常用的 0–1 归一化和“谁好谁坏”的映射可以这么写：

QED

RDKit/TDC 计算出来的 QED 本身就在 ([0,1]) 区间内，所以直接把它当成归一化后的值即可：
[
\text{QED}_{\text{norm}} = \text{QED}
]
值越接近 1 表示“药像度”越好。
SA（Synthetic Accessibility）

RDKit 版本的 SA 分数大致落在 ([1,10]) 区间，数值越大越难合成。
先线性归一化，再取反，让“越好越亮”：
[
\text{SA}{\text{norm}} = 1 - \frac{\text{SA} - \text{SA}{\min}}{\text{SA}{\max} - \text{SA}{\min}}
]
其中 (\text{SA}{\min}\approx 1,\ \text{SA}{\max}\approx 10)。
如果你直接用 1 和 10：(\text{SA}_{\text{norm}} = 1 - \frac{\text{SA} - 1}{9})。
结果在 ([0,1])，越接近 1 表示越容易合成。
Docking Score

Docking 通常越小（越负）越好。假设当前数据的最优值为 (\text{dock}{\min})（最小/最负），最差值为 (\text{dock}{\max})（最大/最正）。
归一化并“颠倒”成越亮越好：
[
\text{Dock}{\text{norm}} = \frac{\text{dock}{\max} - \text{dock}}{\text{dock}{\max} - \text{dock}{\min}}
]
这样最优（(\text{dock} = \text{dock}_{\min})）会映射到 1，最差映射到 0。
实际使用时，只要把当前数据集的 (\text{SA}{\min})、(\text{SA}{\max})、(\text{dock}{\min})、(\text{dock}{\max}) 先统计出来，再带入这些公式，就能生成统一的 0–1 区间颜色映射。